\section{Sovereign Business Map Governance, Entity Linking and Source Lifecycle}
\label{sec:sovereign-business-map-governance}

This increment completes four connected items in Phase~2 of the Sovereign
Business Brain roadmap: \texttt{SBB-0203}, \texttt{SBB-0204},
\texttt{SBB-0206}, and \texttt{SBB-0209}.  It does not create a second business
map.  The existing Boardroom \texttt{BusinessMapContainer} remains the sole
business-map authority, stored inside the customer's Business Container
repository and bound to Boardroom decisions by its exact revision and content
hash.

\subsection{Purpose}

The downloaded ``empty brain'' must be able to learn a company without silently
turning guesses, stale imports, or model-generated associations into company
truth.  The completed layer therefore governs four related objects:

\begin{enumerate}
  \item facts about the business;
  \item relationships between explicit business entities;
  \item approved source passages linked to those entities; and
  \item the retention, legal-hold, and revocation state of each source.
\end{enumerate}

Every material mutation is revision checked.  A stale browser cannot overwrite
a newer map: its request fails with a revision conflict and must reload the
authoritative mother-brain state.  Successful changes produce local hash-bound
governance receipts.

\subsection{Canonical map relationship contract}

The existing \texttt{BusinessMapContainer} now carries a first-class
\texttt{relationships} collection alongside facts, assumptions, unknowns,
conflicts, unanswered fields, and department-discovery gaps.  A governed
relationship records at least:

\begin{itemize}
  \item a stable relationship identifier;
  \item explicit source and destination entity identifiers;
  \item a relationship type;
  \item a source reference and confidence value;
  \item a proposed, approved, or rejected review state;
  \item a verification status and review owner; and
  \item creation and review timestamps.
\end{itemize}

New relationships cannot be proposed without both provenance and confidence.
The owner can approve, correct, reject, or explicitly delete a relationship.
Correction may change an endpoint or the relationship type, but it does not
erase the before-state hash from the receipt ledger.  Deletion requires a
separate positive confirmation.

\subsection{Legacy relationship normalization}

Older Boardroom records may contain useful explicit relationships without the
new governance fields.  A deterministic normalizer upgrades these records in
place.  Missing provenance is labelled as the legacy storage location rather
than misrepresented as an independently verified source.  Its default
confidence is deliberately limited to \(0.5\), and the relationship returns to
the proposed queue for owner review.

The normalizer also reads the existing \texttt{BusinessStructureContainer}.
Where that container already supplies a stable entity identifier, it may create
an explicit structural relationship such as \emph{offers}, \emph{has team},
\emph{uses channel}, or \emph{fulfils through}.  Items without identifiers are
skipped.  No model is asked to infer a missing relationship, and duplicate
source--type--destination triples are not added.

\subsection{Owner-facing map editor}

The management interface is embedded in the existing Boardroom Business Map
panel.  It is not a new dashboard or a browser-local map.  The founder or an
authorised operator can open \emph{Review and edit map} and see plain-language
cards for facts and relationships.  Each card exposes its value, source,
confidence, and review state.

The available controls are:

\begin{itemize}
  \item approve a proposed record;
  \item correct its value or relationship type;
  \item reject it without pretending it disappeared;
  \item delete it after a separate confirmation; and
  \item merge facts with the same field while retaining the duplicate's source
        reference and content hash.
\end{itemize}

The interface also shows the combined review and discovery queue and provides a
bounded \emph{Recover existing relationships} operation.  After every action it
reloads the map from the mother brain.  An owner confirmation used by a Board
meeting becomes stale automatically when the map revision or receipt hash
changes.

\subsection{Deterministic entity linking}

The Business Knowledge index can now link approved document claims to explicit
Boardroom entities.  The entity catalogue is assembled only from Business
Identity and Business Structure records that have stable identifiers and
names.  Supported entity families include the business, services, channels,
revenue streams, costs, payment terms, fulfilment processes, functions, teams,
human agents, and AI agents.

Linking uses case-insensitive exact-name boundaries.  It does not use face
recognition, semantic guessing, or a language model.  A link is created only
when:

\begin{enumerate}
  \item the source claim has already been approved by the owner;
  \item the Boardroom entity has an explicit identifier and name; and
  \item that exact name occurs in the approved claim.
\end{enumerate}

Each link retains the claim identifier, source identifier, source hash, entity
identifier, entity type, deterministic match method, confidence, actor, and
timestamp.  Rebuilding the same links is idempotent.  Pending or revoked claims
cannot become linked business knowledge.

\subsection{Retention, legal hold and source revocation}

Each ingested business source may receive a retention deadline and an explicit
legal-hold flag.  The owner can change that policy through the mother-brain API,
and the change produces a governance receipt.

A source revocation requires exact confirmation.  When authorised, the system:

\begin{enumerate}
  \item marks the source as revoked;
  \item marks every derived claim as revoked and removes it from retrieval;
  \item removes entity links derived from those claims;
  \item deletes the local raw source material only when its resolved path is
        inside that business's bounded raw-source directory; and
  \item retains a non-secret, hash-bound governance receipt.
\end{enumerate}

A legal hold blocks both manual revocation and automatic retention expiry.
The retention worker processes only expired, non-held, non-revoked sources.
Repeated revocation is idempotent.  This separates withdrawal of knowledge from
rewriting historical proof that the governance action occurred.

\subsection{Mother-brain API surface}

The existing Business Twin Data API now exposes bounded operations for fact
merge, relationship proposal and review, legacy normalization, entity-link
rebuild, source-policy updates, confirmed source revocation, and retention
application.  Credentials and provider tokens are not accepted by these
contracts.  The canonical map status also reports its relationship count.

\subsection{Safety and sovereignty properties}

The implementation enforces the following properties:

\begin{itemize}
  \item no parallel business map is introduced;
  \item no legacy relationship is silently promoted to verified truth;
  \item no unapproved document passage participates in entity linking;
  \item no unidentified structure item receives an invented identifier;
  \item no deletion or source revocation occurs without exact confirmation;
  \item no held source is removed by the retention worker;
  \item no raw file path outside the bounded business source directory is
        deleted; and
  \item no browser mutation bypasses optimistic revision control.
\end{itemize}

\subsection{Release and verification status}

This capability is included in AION Sovereign Brain Bootstrap~0.4.0 and Pilot
Fabric~0.56.0.  Automated verification covers legacy normalization, explicit
structure projection, relationship review, duplicate-fact merge, deterministic
entity linking, exclusion of pending claims, source revocation, raw-material
removal, retrieval withdrawal, legal-hold protection, retention expiry, and the
existing Boardroom editor contract.

The next open item in this phase is \texttt{SBB-0207}: one onboarding connector
screen for the already implemented Files, Gmail, Calendar, CRM/Sales, and Xero
read routes, including explicit proof of each connector's map projection.  It is
not claimed complete by this section.
